\documentclass[11pt]{exam}
\input{../Header/packages.sty}
\input{../Header/styleExam.sty}
\input{../Header/mathexpressions.sty}
\input{../Header/headerExam.sty}
%=========================================================
\newcommand{\theme}{Wachstum, Zerfall und Logarithmus}
\newcommand{\teacher}{Jonas Guler}
\newcommand{\class}{3Ga}
\renewcommand{\date}{26.10.26}
\graphicspath{{Algebra/data-algebra},{Geometrie/data-geometrie},{Analysis/data-analysis}} % Setzt den Pfad zu den Bildern fest
%=========================================================
\begin{document}
\fontsize{14}{14}\selectfont
%========================================================
\noindent
\begingroup
\renewcommand{\arraystretch}{1.5}
\begin{tabular*}{\textwidth}{l @{\extracolsep{6pt}}l @{\extracolsep{\fill}}l @{\extracolsep{6pt}}l}
	\multicolumn{2}{l}{{\bf \theme}} & {\bf Note:} & {\answerbox}\\
	\multicolumn{2}{l}{{\bf \teacher}} & & \\
	{\bf Klasse:} & {\class} &  {\bf Name:} & {\rule{5cm}{0.5pt}}\\
	{\bf Dauer:} & $45$ Minuten & {\bf Datum:} & {\date}\\[8pt]
\end{tabular*}
\endgroup
\rule[1ex]{\textwidth}{1pt}
%=========================================================
\begin{center}
    \textbf{Bewertungstabelle}\\
    \small Wird von der Lehrperson ausgefüllt\\
    \vspace{5pt}
    \addpoints
    \gradetable[h][questions]\\
    \vspace{5pt}
    Der Rechenweg muss \textbf{vollständig und nachvollziehbar} sein.\\
    Lösen Sie die Prüfung mit einen schwarzen oder blauen Kugelschreiber.\\
    Als Hilfsmittel ist der Taschenrechner (TI - Nspire CX-II CAS) erlaubt.
    %Als Hilfsmittel ist der kleine Taschenrechner (TI-30X o.ä.) erlaubt.
\end{center}
\rule[1ex]{\textwidth}{1pt}

\vspace{10cm}
\begin{center}
    %\Huge{Repetitionsprüfung}
\end{center}

\newpage
\begin{questions}
%=========================================================
%---------------------TESTFRAGEN
%=========================================================
\noaddpoints
\question[\totalpoints]
\addpoints
Beantworte für die beiden Funktionen $f:y=a\*x$ und $g:y=a^x$ die folgenden Fragen: Wie verändert sich der Funktionswert, ...
\begin{parts}
	\part[2] ... wenn man den $x-$Wert um 4 verringert?
	\part[2] ... wenn man den $x-$Wert verdreifacht?
\end{parts}
\antwort{8.5}

\question[3]
Im folgenden sind drei Wertetabellen gegeben. Entscheide für jede Wertetabelle, ob es sich um eine Exponentialfunktion handelt. Wenn ja, gib die Funktionsgleichung an.
\begin{parcols}[3]
	\part \
	\begin{table}[H]
		\begin{tabular}{|c|c|c|c|c|}\hline
			$x$ & 1 & 2 & 3 & 4 \\\hline
			$f(x)$ & 3 & 9 & 27 & 81 \\\hline
		\end{tabular}
	\end{table}
	\part \
	\begin{table}[H]
		\begin{tabular}{|c|c|c|c|c|}\hline
			$x$ & 1 & 2 & 3 & 4 \\\hline
			$f(x)$ & -7 & -2 & 3 & 8 \\\hline
		\end{tabular}
	\end{table}
	\part\
	\begin{table}[H]
		\begin{tabular}{|c|c|c|c|c|}\hline
			$x$ & 1 & 2 & 3 & 4 \\\hline
			$f(x)$ & 32 & 8 & 2 & $\frac{2}{4}$ \\\hline
		\end{tabular}
	\end{table}
\end{parcols}
\antwort{7.5}


%-----------------------------------------------------------------------------
\newpage
\question[5]
Ein neu entwickeltes Schmerzmittel namens Analgetin wird vom Körper so abgebaut, dass in $5$ Stunden jeweils $20\%$ der noch im Blut vorhandenen Wirkstoffmenge abgebaut werden.
\begin{parts}
	\part Berechne die Halbwertszeit des Wirkstoffs im Blut.
	\part Wie viel Prozent der noch vorhandenen Wirkstoffmenge werden im Lauf des nächsten Tages ($24$ Stunden) abgebaut?
\end{parts}
\antwort{20.5}


%-----------------------------------------------------------------------------
\newpage
\question[6]
Die Verbreitung eines viralen Videos auf einer Online-Plattform lässt sich in der Anfangsphase durch eine Exponentialfunktion beschreiben. $2$ Tage nach dem Hochladen wurde das Video $40'000$-mal aufgerufen. $2$ Tage später sind es bereits $900'000$ Aufrufe.
\begin{parts}
	\part Berechne die Anzahl der Aufrufe $5$ Tage nach dem Hochladen.
	\part In welcher Zeit verdreifacht sich die Anzahl der Aufrufe?
	\part Um wie viele Prozent wächst die Anzahl der Aufrufe innerhalb eines Tages?
\end{parts}
\antwort{19.5}


%-----------------------------------------------------------------------------
\newpage
\question[4]
Als eine Gemeinde ihr Förderprogramm für Solaranlagen startete, waren auf ihren Hausdächern $250$ Solaranlagen installiert. Ein Jahr später waren es bereits $310$.
\begin{parts}
	\part Wie viele Solaranlagen werden es am Tag des $10-$jährigen Jubiläums des Förderprogramms sein, wenn wir ein lineares Wachstum annehmen?
	\part Wie viele Solaranlagen werden es am Jubiläumstag sein, wenn wir ein exponentielles Wachstum annehmen?
\end{parts}
\antwort{8}

\question[4]
In das leere Feld ist jeweils ein Ausdruck so einzufüllen, dass insgesamt eine gültige Aussage entsteht.\\[-10pt]
\begin{parts}
	\part $\ds\log{a}{\gap{3}}=3\*\log{a}{u}-2\*\log{a}{v}-5\*\log{a}{w}$\\
	\part $\ds\log{a}{\sqrt[5]{v^3}}=\gap{3}\*\log{a}{v}$\\
	\part $\ds\log{a}{a^{-\frac{7}{11}}}=\gap{3}$\\
	\part $\ds\gap{3}=\frac{2\*\ln{y}}{\ln{5}}$\\
\end{parts}
\antwort{4.5}


%-----------------------------------------------------------------------------
\newpage
\noaddpoints
\question[\totalpoints]
\addpoints
{\bf Bewertung:} Diese Aufgabe ist ohne Taschenrechner zu lösen.
\begin{parts}
	\part[2] Schreibe den folgenden Logarithmus in möglichst viele Summanden aus.
	\[\log{b}{\frac{b^5\*\sqrt{c^5}}{b\*\sqrt[3]{c}}}\]
\end{parts}
\antwort{9}
\begin{parts}
	\setcounter{partno}{1}\vspace{5pt}
	\part[2] Fasse den Term zu einem Logarithmus zusammen und vereinfache so weit wie möglich. \[\ln{x+2}+2\*\ln{x-2}-\ln{x+2}-2\*\ln{x-2}\]
\end{parts}
\antwort{9}


%-----------------------------------------------------------------------------
\newpage
\question[4]
{\bf Bewertung:} Diese Aufgabe ist ohne Taschenrechner zu lösen.\\
Löse die folgende Gleichung für eine beliebige Basis $a$ nach $x$ auf. \[\log{a}{x^2-5}+\log{a}{x+5}=\log{a}{x+5}\] {\textit{Tipp: Forme die Gleichung sinnvoll mithilfe der Logarithmusgesetzen um.}}
\antwort{21.5}

\end{questions}
\end{document}